% Standalone manuscript section. Compile with pdfLaTeX.
% To insert into an existing paper, copy from \section through the appendix,
% adapting bibliography keys and removing the standalone preamble.
\documentclass[11pt]{article}
\usepackage[a4paper,margin=25mm]{geometry}
\usepackage{amsmath,amssymb,bm}
\usepackage[hidelinks]{hyperref}
\newcommand{\Fa}{\mathbf F_a}
\newcommand{\Fd}{\mathbf F_d}
\newcommand{\pos}{\mathbf p}
\newcommand{\dd}{\mathrm d}
\title{Single-User, Single-Target Robust Hybrid ISAC\\with Fluid Transmit Antennas}
\author{}
\date{}
\begin{document}
\maketitle
\noindent\textbf{Scope.} This proposed model retains the robust hybrid-beamforming
formulation of \cite{wang2026}, specializes it to one user and one target, and
replaces only the fixed transmit ULA coordinates by adjustable one-dimensional
fluid-antenna coordinates. The radar receive array remains a fixed ULA.
Continuous spatial reconfiguration is motivated by \cite{zhou2024}; we do not
adopt that paper's different sensing objective or fluid user antenna.

\section{System Model}
\subsection{Network and antenna geometry}
Consider a downlink integrated sensing and communication (ISAC) system with a
dual-function base station (BS), one fixed single-antenna communication user,
one point-like target, and $J\geq0$ clutter sources. The BS employs $N_t$ fluid
transmit antennas, $L$ RF chains, and $N_r$ fixed radar receive antennas.
We consider narrowband far-field operation. Antenna locations are configured
before a transmission block and held fixed during that block. The array
aperture is small relative to propagation distances, so path angles and gains
are held fixed as candidate transmit positions are evaluated.

Let $\pos=[p_1,\ldots,p_{N_t}]^T\in\mathbb R^{N_t}$ collect the transmit
coordinates on an aperture of length $A$. Label the antennas in spatial order
and define
\begin{equation}
 \mathcal P=\left\{\pos:\ 0\leq p_1,\quad
 p_{n+1}-p_n\geq d_{\min},\ n=1,\ldots,N_t-1,\quad p_{N_t}\leq A\right\}.
 \label{eq:positions}
\end{equation}
Here $d_{\min}>0$ is the minimum antenna separation. The necessary and
sufficient condition for this geometric set to be nonempty is
$A\geq(N_t-1)d_{\min}$. This ordered one-dimensional set is convex;
the complete optimization problem will nevertheless be nonconvex.
We assume accurate position control and neglect mutual coupling in this
initial model. The receive ULA has spacing $d_r$.

\subsection{Position-dependent transmit response}
Let $\lambda$ be the wavelength, $k_0=2\pi/\lambda$, and $\theta$ denote an
angle measured from array broadside. Following the phase convention in
\cite{wang2026}, the nominal transmit steering vector is
\begin{equation}
 \widehat{\mathbf a}_t(\theta;\pos)
 =\left[e^{-jk_0p_1\sin\theta},\ldots,
 e^{-jk_0p_{N_t}\sin\theta}\right]^T.
 \label{eq:txsteering}
\end{equation}
Indeed, the relative path-length term at antenna $n$ is $p_n\sin\theta$;
multiplication by $k_0$ gives its relative phase. For the fixed receive ULA,
\begin{equation}
 \widehat{\mathbf a}_r(\theta)
 =\left[1,e^{-jk_0d_r\sin\theta},\ldots,
 e^{-jk_0(N_r-1)d_r\sin\theta}\right]^T.
 \label{eq:rxsteering}
\end{equation}
Setting $p_n=(n-1)d$ in \eqref{eq:txsteering} recovers the original transmit
ULA with spacing $d$. Thus the fixed-array geometry is a special case whenever
it satisfies \eqref{eq:positions}.

\subsection{Position-dependent communication channel}
To evaluate different geometries consistently, model the nominal communication
channel through $K_c$ propagation paths:
\begin{equation}
 \widehat{\mathbf h}^{H}(\pos)
 =\sum_{k=1}^{K_c}\beta_k
 \widehat{\mathbf a}_t^T(\vartheta_k;\pos),
 \qquad \widehat{\mathbf h}(\pos)\in\mathbb C^{N_t}.
 \label{eq:commchannel}
\end{equation}
Here $\beta_k$ includes the complex path gain at the aperture reference and
$\vartheta_k$ is its departure angle. Consequently,
$[\widehat{\mathbf h}^{H}(\pos)]_n=
\sum_k\beta_ke^{-jk_0p_n\sin\vartheta_k}$.
The gains and angles are fixed within each scenario; channels must not be
independently redrawn after position updates. The actual channel is
\begin{equation}
 \mathbf h(\pos)=\widehat{\mathbf h}(\pos)+\bm\delta_h.
 \label{eq:herror}
\end{equation}
This spatial construction is the channel-model addition required to make
fluid positioning meaningful; the received-signal structure remains unchanged.

\subsection{Hybrid transmit signal}
Let $z_c$ and $z_s$ denote the zero-mean communication symbol and dedicated
sensing symbol. Assume
\begin{equation}
 \mathbb E|z_c|^2=\mathbb E|z_s|^2=1,
 \qquad \mathbb E[z_cz_s^*]=0.
\end{equation}
The analog beamformer is $\Fa\in\mathbb C^{N_t\times L}$ and the digital
beamformer is $\Fd=[\mathbf f_c,\mathbf f_s]\in\mathbb C^{L\times2}$.
For the initial two-stream implementation, choose $2\leq L\leq N_t$.
The transmitted signal is
\begin{equation}
 \mathbf x=\Fa\mathbf f_cz_c+\Fa\mathbf f_sz_s
 =\Fa\Fd\mathbf z,\qquad \mathbf z=[z_c,z_s]^T.
 \label{eq:transmit}
\end{equation}
Writing $\mathbf v_c=\Fa\mathbf f_c$ and $\mathbf v_s=\Fa\mathbf f_s$,
the transmit covariance and average radiated power follow as
\begin{align}
 \mathbf R_x&=\mathbb E[\mathbf x\mathbf x^H]
 =\mathbf v_c\mathbf v_c^H+\mathbf v_s\mathbf v_s^H
 =\Fa\Fd\Fd^H\Fa^H,\label{eq:covariance}\\
 \mathbb E\|\mathbf x\|_2^2
 &=\operatorname{Tr}(\mathbf R_x)=\|\Fa\Fd\|_F^2.
 \label{eq:power}
\end{align}
The cross terms vanish by the uncorrelated-stream assumption. We retain the
original constraints $|[\Fa]_{n,\ell}|=1$ and
$\|\Fa\Fd\|_F^2=P_{\mathrm{BS}}$.

\subsection{Communication received signal and rate}
With noise $n_c\sim\mathcal{CN}(0,\sigma_c^2)$, the user receives
\begin{equation}
 y_c=\mathbf h^H(\pos)\mathbf x+n_c
 =\underbrace{\mathbf h^H(\pos)\Fa\mathbf f_cz_c}_{\text{desired signal}}
 +\underbrace{\mathbf h^H(\pos)\Fa\mathbf f_sz_s}_{\text{sensing interference}}
 +n_c.
 \label{eq:commreceived}
\end{equation}
The independent unit-power streams give desired and interfering powers
$|\mathbf h^H\Fa\mathbf f_c|^2$ and $|\mathbf h^H\Fa\mathbf f_s|^2$,
respectively. Thus,
\begin{equation}
 \mathrm{SINR}_c=
 \frac{|\mathbf h^H(\pos)\Fa\mathbf f_c|^2}
 {|\mathbf h^H(\pos)\Fa\mathbf f_s|^2+\sigma_c^2},
 \qquad \gamma_c=\log(1+\mathrm{SINR}_c).
 \label{eq:commsinr}
\end{equation}
There is no inter-user interference. As in the reference model, the user treats
the dedicated sensing stream as interference. We use natural logarithms
throughout, giving the communication metric in nats/s/Hz; replacing both
communication and sensing logarithms by $\log_2$ scales both metrics by the
same constant. The standard achievable-rate interpretation assumes suitable
Gaussian signaling and interference treatment.

\subsection{Radar echoes and sensing metric}
Index the target by $l=0$ and clutter sources by $l=1,\ldots,J$. Retain the
bounded steering-response model of \cite{wang2026}:
\begin{align}
 \mathbf a_{t,l}(\pos)&=
 \widehat{\mathbf a}_t(\theta_l;\pos)+\bm\delta_{t,l},\\
 \mathbf a_{r,l}&=\widehat{\mathbf a}_r(\theta_l)+\bm\delta_{r,l}.
\end{align}
The round-trip channel becomes
\begin{equation}
 \mathbf H_l(\pos)=\alpha_l\mathbf a_{r,l}\mathbf a_{t,l}^T(\pos)
 \in\mathbb C^{N_r\times N_t},\qquad l=0,\ldots,J,
 \label{eq:radarchannel}
\end{equation}
where $\alpha_l$ contains the reflection and round-trip propagation gain.
The receive ULA and the transpose convention are unchanged from the reference.
The echo vector and filter output are
\begin{align}
 \mathbf y_r&=\mathbf H_0(\pos)\mathbf x
 +\sum_{j=1}^{J}\mathbf H_j(\pos)\mathbf x+\mathbf n_r,\\
 y_s&=\mathbf w^H\mathbf y_r,\qquad
 \mathbf n_r\sim\mathcal{CN}(\mathbf0,\sigma_r^2\mathbf I_{N_r}),
 \quad \|\mathbf w\|_2^2=1.
\end{align}
The effective echo model neglects residual transmitter self-interference and
does not explicitly resolve delay or Doppler. The desired target power is
\begin{align}
 P_0&=\mathbb E|\mathbf w^H\mathbf H_0\mathbf x|^2
 =\mathbf w^H\mathbf H_0\mathbf R_x\mathbf H_0^H\mathbf w\\
 &=\|\mathbf w^H\mathbf H_0\Fa\Fd\|_2^2
 =|\mathbf w^H\mathbf H_0\mathbf v_c|^2
  +|\mathbf w^H\mathbf H_0\mathbf v_s|^2.
\end{align}
Both streams therefore illuminate the target. Retaining the reference's
sum-of-powers treatment of effective clutter contributions gives
\begin{equation}
 \mathrm{SINR}_s=
 \frac{\|\mathbf w^H\mathbf H_0(\pos)\Fa\Fd\|_2^2}
 {\sum_{j=1}^{J}\|\mathbf w^H\mathbf H_j(\pos)\Fa\Fd\|_2^2
  +\sigma_r^2\|\mathbf w\|_2^2},
 \qquad \gamma_s=\log(1+\mathrm{SINR}_s).
 \label{eq:sensingsinr}
\end{equation}
Here $\gamma_s$ is the inherited MI-based sensing metric, rather than a new
claim about exact parameter-estimation mutual information. The expression
retains the reference approximation that effective echo cross-correlations
are negligible. Independent source streams alone do not guarantee this:
coherent clutter would additionally contribute terms
$\mathbf w^H\mathbf H_j\mathbf R_x\mathbf H_k^H\mathbf w$ for $j\ne k$.
We preserve the original metric to isolate the effect of fluid positioning.

\subsection{Uncertainty sets}
Keep the uncertainty balls unchanged in form:
\begin{align}
 \mathcal A_h&=\{\bm\delta_h\in\mathbb C^{N_t}:
 \|\bm\delta_h\|_2\leq\epsilon_h\sqrt{N_t}\},\\
 \mathcal A_t&=\prod_{l=0}^{J}\{\bm\delta_{t,l}\in\mathbb C^{N_t}:
 \|\bm\delta_{t,l}\|_2\leq\epsilon_t\sqrt{N_t}\},\\
 \mathcal A_r&=\prod_{l=0}^{J}\{\bm\delta_{r,l}\in\mathbb C^{N_r}:
 \|\bm\delta_{r,l}\|_2\leq\epsilon_r\sqrt{N_r}\}.
\end{align}
Write $\bm\Delta_t=\{\bm\delta_{t,l}\}_{l=0}^{J}$ and
$\bm\Delta_r=\{\bm\delta_{r,l}\}_{l=0}^{J}$.
The radii are specified model parameters in the initial comparison.
They are kept identical for fixed and fluid arrays. These are unstructured
response-error sets; their elements need not correspond to a single physical
angle error. No separate antenna-position error is introduced at this stage.

\subsection{Joint robust optimization}
For a prescribed tradeoff $0<\rho<1$, the proposed extension is
\begin{equation}
\begin{aligned}
 \underset{\pos,\Fa,\Fd,\mathbf w}{\operatorname{maximize}}\quad
 &\rho\min_{\bm\delta_h\in\mathcal A_h}\gamma_c
 +(1-\rho)\min_{\substack{\bm\Delta_t\in\mathcal A_t\\
                          \bm\Delta_r\in\mathcal A_r}}\gamma_s\\
 \text{subject to}\quad
 &|[\Fa]_{n,\ell}|=1,\quad\forall n,\ell,\\
 &\|\Fa\Fd\|_F^2=P_{\mathrm{BS}},\qquad\|\mathbf w\|_2^2=1,\\
 &\pos\in\mathcal P.
\end{aligned}
\label{eq:problem}
\end{equation}
Relative to the single-user, single-target specialization of \cite{wang2026},
the only new design block is $\pos$. The received-signal formulas, uncertainty
sets, hybrid constraints, and weighted robust objective retain their original
structure. Position dependence enters through \eqref{eq:commchannel} and
\eqref{eq:radarchannel}. Fixing $p_n=(n-1)d$ recovers the fixed-array geometry
under the same spatial channel construction. This does not by itself reproduce
an unrelated position-independent channel-generation distribution.

\appendix
\section{Additional Derivatives Introduced by Fluid Positions}
The following derivatives hold for fixed beamformers and fixed additive
uncertainty realizations. Let $\mathbf e_n$ be the $n$th coordinate vector
and let a prime denote differentiation with respect to real coordinate $p_n$.
From \eqref{eq:txsteering},
\begin{equation}
 \mathbf a'_{t,l,n}
 =-jk_0\sin\theta_l\,
 [\widehat{\mathbf a}_t(\theta_l;\pos)]_n\mathbf e_n.
\end{equation}
Only the nominal response is differentiated, because the additive errors
are held fixed in this partial derivative. Similarly,
\begin{align}
 (\mathbf h^H)'_n
 &=\sum_{k=1}^{K_c}\beta_k(-jk_0\sin\vartheta_k)
 e^{-jk_0p_n\sin\vartheta_k}\mathbf e_n^T,\\
 \mathbf H'_{l,n}
 &=\alpha_l\mathbf a_{r,l}(\mathbf a'_{t,l,n})^T.
\end{align}
For any differentiable complex scalar $b(p_n)$,
\begin{equation}
 \frac{\partial |b|^2}{\partial p_n}
 =2\operatorname{Re}\{b^*b'\}.
\end{equation}
Define $b_c=\mathbf h^H\mathbf v_c$, $b_s=\mathbf h^H\mathbf v_s$,
$N_c=|b_c|^2$, and $D_c=|b_s|^2+\sigma_c^2$. Then
\begin{align}
 N'_{c,n}&=2\operatorname{Re}\{b_c^*(\mathbf h^H)'_n\mathbf v_c\},\\
 D'_{c,n}&=2\operatorname{Re}\{b_s^*(\mathbf h^H)'_n\mathbf v_s\},\\
 \frac{\partial\gamma_c}{\partial p_n}
 &=\frac{D_cN'_{c,n}-N_cD'_{c,n}}{D_c(D_c+N_c)}.
\end{align}
For sensing, put $\mathbf q_l=\mathbf w^H\mathbf H_l\Fa\Fd$ and
$\mathbf q'_{l,n}=\mathbf w^H\mathbf H'_{l,n}\Fa\Fd$. Define
$N_s=\|\mathbf q_0\|_2^2$ and
$D_s=\sum_{j=1}^{J}\|\mathbf q_j\|_2^2+\sigma_r^2\|\mathbf w\|_2^2$.
It follows that
\begin{align}
 N'_{s,n}&=2\operatorname{Re}\{\mathbf q'_{0,n}\mathbf q_0^H\},\\
 D'_{s,n}&=2\operatorname{Re}\left\{\sum_{j=1}^{J}
 \mathbf q'_{j,n}\mathbf q_j^H\right\},\\
 \frac{\partial\gamma_s}{\partial p_n}
 &=\frac{D_sN'_{s,n}-N_sD'_{s,n}}{D_s(D_s+N_s)}.
\end{align}
Thus the fixed-error outer utility gradient is
\begin{equation}
 \frac{\partial\gamma}{\partial p_n}
 =\rho\frac{\partial\gamma_c}{\partial p_n}
 +(1-\rho)\frac{\partial\gamma_s}{\partial p_n}.
\end{equation}
For base-two logarithms, divide these utility derivatives by $\ln 2$.
Using an approximately minimized error realization in these expressions gives
the position-update direction of an alternating PGDA scheme; it is not an
automatic certificate of the exact derivative of the globally minimized
robust objective. The inherited beamforming/error gradients keep their
algebraic structure, evaluated with the position-dependent channels.

\begin{thebibliography}{2}
\bibitem{wang2026}
L. Wang, S. A. Vorobyov, and E. Ollila,
``L2O robust hybrid beamforming for ISAC,''
\emph{IEEE Transactions on Wireless Communications}, 2026,
doi: \href{https://doi.org/10.1109/TWC.2026.3702925}{10.1109/TWC.2026.3702925}.
\bibitem{zhou2024}
L. Zhou, J. Yao, M. Jin, T. Wu, and K.-K. Wong,
``Fluid antenna-assisted ISAC systems,''
\emph{IEEE Wireless Communications Letters}, vol. 13, no. 12, 2024,
doi: \href{https://doi.org/10.1109/LWC.2024.3476148}{10.1109/LWC.2024.3476148}.
\end{thebibliography}
\end{document}